\documentclass[fontsize=15pt]{jlreq}
%プリアンブル
\usepackage{amsmath}
\usepackage{mathtools}
\pagestyle{empty}
\usepackage[margin=10mm]{geometry}
\usepackage{tikz}
\usetikzlibrary{arrows.meta, calc}
\usepackage{here}
\usepackage{tabularx}

\newcommand{\m}[1]{\raisebox{0.1em}{\textcircled{\scriptsize #1}}}



%本文
\begin{document}
\begin{center}
{\LARGE {振り返り}} \\
-科学技術と人間- %単元ごとに変えるか消す
\end{center}
\vspace{1em}
\begin{flushright}
名前：\underline{\hspace{5cm}}
\end{flushright}

% 問題部分
\begin{enumerate}
  \item プラスチックの主な原料をひとつかけ．\\
  (\hspace{3cm})
  
  \vspace{2em}

  \item 7つあるプラスチックの性質のうち，書けるだけかけ．
  (一応7個分の回答欄を用意しておく．最低3つは書いて)\\
  \begin{tabularx}{\textwidth}{XX}
    \m{1}&\m{2}\\
    \m{3}&\m{4}\\
    \m{5}&\m{6}\\
    \m{7}&\\
  \end{tabularx}

  \vspace{2em}

  \item 放射線の性質と，放射線を出す能力の名称を書け．\\
  性質：  \\


  \noindent 名称：

  \vspace{2em}

  \item 放射線についての記述のうち，間違っているものを1つ選べ．
  \begin{itemize}
    \item[ア：] 受けた放射線の人体に対する影響を表す単位の1つに
    [Sv](シーベルト)がある．
    \item[イ：] 家では，放射線の影響を受けることはない．
    \item[ウ：] どのような物質であっても，放射線の影響を受けないということはない．   
  \end{itemize}
  (\hspace{3em})

  \vspace{2em}

  \item 再生可能エネルギーを書けるだけ書け．
  
\end{enumerate}

\end{document}